\documentclass[11pt,a4paper]{article}
\usepackage[T1]{fontenc}
\usepackage[utf8]{inputenc}
\usepackage{lmodern}
\usepackage{amsmath,amssymb,amsthm,mathtools,mathrsfs}
\usepackage{graphicx,float,xcolor,multicol}
\usepackage[shortlabels]{enumitem}
\usepackage[margin=27mm]{geometry}
\usepackage{microtype,needspace,placeins}
\usepackage{tikz,tikz-cd}
\usetikzlibrary{arrows.meta,calc,positioning,patterns,decorations.pathreplacing}
\usepackage[colorlinks=true,linkcolor=teal,urlcolor=teal,citecolor=teal]{hyperref}
\setlength{\parindent}{0pt}
\setlength{\parskip}{.6em}
\setcounter{secnumdepth}{0}
\allowdisplaybreaks
\raggedbottom
\providecommand{\cross}{\times}
\DeclareUnicodeCharacter{2020}{\ensuremath{\dagger}}
\DeclareMathOperator{\supp}{supp}
\DeclareMathOperator*{\esssup}{ess\,sup}
\providecommand{\chapter}{\section}
\newenvironment{tool}[2]{\subsection*{Tool #1 --- #2}}{}
\newcommand{\ToolRef}[1]{Tool #1}
\newcommand{\FigureTag}[2]{\textit{Figure #1}}
\newcommand{\Result}[1]{\subsection*{#1}}
\newcommand{\Application}[1]{\subsubsection*{Application\if\relax\detokenize{#1}\relax\else\space--- #1\fi}}
\newcommand{\ProofHeading}{\par\textit{Proof.}\space}
\newcommand{\ProofOf}[1]{\par\textit{Proof of #1.}\space}
\newcommand{\RemarkHeading}{\par\textbf{Remark.}\space}
\newcommand{\NoteHeading}{\par\textbf{Note.}\space}
\newcommand{\ExerciseHeading}{\par\textbf{Exercise.}\space}

% Rendering support only; mathematical text is in each independent post.
\usepackage{booktabs,array,longtable}
\definecolor{HandoutBlue}{HTML}{2F6690}
\definecolor{HandoutNavy}{HTML}{17365D}
\newtheorem{theorem}{Theorem}
\newtheorem{lemma}[theorem]{Lemma}
\newtheorem{proposition}[theorem]{Proposition}
\newtheorem{corollary}[theorem]{Corollary}
\newtheorem{conjecture}[theorem]{Conjecture}
\newtheorem{definition}{Definition}
\newtheorem{example}{Example}
\newtheorem{namedtoolcounter}{Tool}
\newenvironment{namedtool}[1]{\begin{namedtoolcounter}[#1]}{\end{namedtoolcounter}}
\newtheorem*{claim}{Claim}
\newtheorem*{remark}{Remark}
\newtheorem*{exercise}{Exercise}
\newtheorem*{question}{Question}
\newtheorem*{observation}{Observation}
\newcommand{\SourceChapter}[2]{\section*{#2}}
\newcommand{\Lecture}[2]{\section*{Lecture #1: #2}}
\newcommand{\unclear}[1]{\textit{[unclear: #1]}}
\newcommand{\sourcegap}{\textit{[unfinished in the source]}}
\DeclareMathOperator{\dist}{dist}
\DeclareMathOperator{\id}{id}
\DeclareMathOperator{\Hol}{Hol}
\DeclareMathOperator{\Per}{Per}
\DeclareMathOperator{\Fix}{Fix}
\DeclareMathOperator{\diam}{diam}
\DeclareMathOperator{\Var}{Var}
\DeclareMathOperator{\Leb}{Leb}
\DeclareMathOperator{\Hom}{Hom}
\DeclareMathOperator{\End}{End}
\DeclareMathOperator{\Ann}{Ann}
\DeclareMathOperator{\Spec}{Spec}
\DeclareMathOperator{\Max}{Max}
\DeclareMathOperator{\Ob}{Ob}
\DeclareMathOperator{\Tor}{Tor}
\DeclareMathOperator{\Ass}{Ass}
\DeclareMathOperator{\Mor}{Mor}
\DeclareMathOperator{\Fun}{Fun}
\DeclareMathOperator{\Nat}{Nat}
\DeclareMathOperator{\Cone}{Cone}
\DeclareMathOperator{\MCG}{MCG}
\DeclareMathOperator{\Aut}{Aut}
\DeclareMathOperator{\Deck}{Deck}
\DeclareMathOperator{\SL}{SL}
\DeclareMathOperator{\PSL}{PSL}
\DeclareMathOperator{\Area}{Area}
\DeclareMathOperator{\im}{im}
\DeclareMathOperator{\PSh}{PSh}
\newcommand{\CRing}{\mathsf{CRings}}
\newcommand{\AMod}{A\text{-}\mathsf{Mod}}
\newcommand{\Ab}{\mathsf{Ab}}
\newcommand{\Z}{\mathbb Z}
\newcommand{\Q}{\mathbb Q}
\newcommand{\R}{\mathbb R}
\newcommand{\C}{\mathbb C}
\newcommand{\N}{\mathbb N}
\newcommand{\kk}{\mathbb k}
\renewcommand{\aa}{\mathfrak a}
\newcommand{\bb}{\mathfrak b}
\newcommand{\pp}{\mathfrak p}
\newcommand{\qq}{\mathfrak q}
\newcommand{\mm}{\mathfrak m}
\newcommand{\nn}{\mathfrak n}
\newcommand{\rad}{\mathrm r}
\newcommand{\cat}[1]{\mathsf{#1}}
\setcounter{secnumdepth}{0}
\allowdisplaybreaks

\graphicspath{{figures/lemma-book/}}
\begin{document}
\section*{Elementary Number Theory}
% BLOG-CONTENT-BEGIN
\subsection{Elementary number theory}
\paragraph{Application: AMTI, inter-final round (2014).}
If $\gcd(a+b,a-b)=1$, find
\[
\gcd\bigl(2a+(2a+1)(a^2-b^2),\;2a(a^2-b^2+2a)(a^2-b^2)\bigr).
\]
We have
\begin{align*}
\gcd(a+b,a-b)=1
&\Longrightarrow\gcd\bigl(2a,(a+b)(a-b)\bigr)=1\\
&\Longrightarrow\gcd\bigl(2a+a^2-b^2,2a(a^2-b^2)\bigr)=1\\
&\Longrightarrow\gcd\bigl(2a+a^2-b^2+2a(a^2-b^2),
 2a(a^2-b^2+2a)(a^2-b^2)\bigr)=1.
\end{align*}
Thus the required GCD is $1$.

\paragraph{Application: British Mathematical Olympiad, round 2 (2009).}
Find all solutions in nonnegative integers $(a,b)$ of
$\sqrt a+\sqrt b=\sqrt{2009}$.
Since $2009=49\cdot41=7^2\cdot41$, we have $\sqrt{2009}=7\sqrt{41}$.
Split seven as a sum of two nonnegative integers to obtain
\[
(a,b)\in\{(656,369),(1025,164),(1476,41),(2009,0),
(369,656),(164,1025),(41,1476),(0,2009)\}.
\]
\emph{Quite elegantly done.}

\setcounter{theorem}{2}\begin{lemma}[Consecutive composite integers]
$n!+2,n!+3,n!+4,\ldots,n!+n$ are $n-1$ consecutive composite integers.
\end{lemma}
No special use or application.

\setcounter{theorem}{3}\begin{lemma}[A small prime factor]
The smallest prime factor of a nonprime $n$ is at most $\sqrt n$.
\end{lemma}
% Source page 3 - left
The fourth lemma can save some time when factorizing a number.

\setcounter{theorem}{4}\begin{lemma}[Primes modulo six]
All primes $p>3$ have the form $6n\pm1$.
\end{lemma}
\paragraph{Application: RMO (1998).}
Let $p_1,p_2,\ldots,p_n$ be $n$ primes, all greater than $5$, such that
$6$ divides $p_1^2+\cdots+p_n^2$. Prove that $6\mid n$.
For every $i$, $p_i^2\equiv1\pmod6$. Hence
$\sum_{i=1}^n p_i^2\equiv n\pmod6$, and the divisibility assumption gives $n\equiv0\pmod6$.

\setcounter{theorem}{5}\begin{lemma}[Fermat's congruence]
Let $a$ be a positive integer and $p$ a prime. Then $a^p\equiv a\pmod p$.
If $\gcd(a,p)=1$, then $a^{p-1}\equiv1\pmod p$.
\end{lemma}

\setcounter{theorem}{6}\begin{lemma}[Wilson's theorem]
If $p$ is prime, then $(p-1)!\equiv-1\pmod p$.
\end{lemma}

% Source page 3 - right
\setcounter{theorem}{7}\begin{lemma}[Euler's totient]
Euler's $\varphi$-function counts the numbers less than $m$ and relatively prime to it.
If $m=\prod_{i=1}^n p_i^{k_i}$, then
\[
\varphi(m)=m\left(1-\frac1{p_1}\right)\cdots\left(1-\frac1{p_n}\right).
\]
If $(a,m)=1$, then $a^{\varphi(m)}\equiv1\pmod m$.
\end{lemma}
This has quite a lot of applications.


\setcounter{theorem}{9}\begin{lemma}[Sophie Germain's identity]
$a^4+4b^4=(a^2+2b^2+2ab)(a^2+2b^2-2ab)$.
\end{lemma}

\setcounter{theorem}{10}\begin{lemma}[A polynomial factorization]
\begin{align*}
x^5+x^4+1
&=\bigl((x+1)^3-3(x+1)^2+2(x+1)+1\bigr)
  \bigl((x+1)^2-(x+1)+1\bigr)\\
&=(x^2+x+1)(x^3-x+1).
\end{align*}
\end{lemma}
\paragraph{Application.}
Show that $n(n-1)^4+1$ is composite for $n>2$.
\[
n(n-1)^4+1=(n-1)^5+(n-1)^4+1
=(n^3-3n^2+2n+1)(n^2-n+1).
\]
Since $n^3-3n^2+2n=n(n-1)(n-2)>1$ and $n^2-n=n(n-1)>1$, the result follows.


% Source page 4 - left
\setcounter{theorem}{12}\begin{lemma}[A cubic identity]
$(a+b+c)^3=a^3+b^3+c^3+3(a+b)(b+c)(c+a)$.
\end{lemma}
\paragraph{Application.}
If $6\mid a+b+c$, then $6\mid a^3+b^3+c^3$.
Indeed, divisibility of $a^3+b^3+c^3$ by $6$ is equivalent to divisibility of
$(a+b+c)^3-3(a+b)(b+c)(c+a)$, hence to that of $3(a+b)(b+c)(c+a)$.
By the box principle, two of $a,b,c$ have the same parity, so the latter product is divisible by $6$.

\subsection{A beautiful application of modular arithmetic}
Prove that $641\mid2^{32}+1$.
Since $640\equiv-1\pmod{641}$,
\[
5\cdot2^7\equiv-1\pmod{641}
\Longrightarrow(5\cdot2^7)^4\equiv1\pmod{641}
\Longrightarrow625\cdot2^{28}\equiv1\pmod{641}.
\]
But $625\equiv-16\pmod{641}$, so
$625\cdot2^{28}\equiv-16\cdot2^{28}=-2^{32}\equiv1\pmod{641}$.
Consequently $2^{32}+1\equiv0\pmod{641}$.



\setcounter{theorem}{15}\begin{lemma}[Squarefree factor]
If $n^2=ab$, where $a<b$ and $a$ is not divisible by a perfect square greater than $1$, then $a\mid b$.
\end{lemma}
\paragraph{Application.}
The sum of the squares of five successive positive integers is not a square.
\[
(n-2)^2+(n-1)^2+n^2+(n+1)^2+(n+2)^2=5(n^2+2).
\]
If this is a square, $5\mid n^2+2$, so $n^2\equiv-2\pmod5$, a contradiction.

\setcounter{theorem}{16}\begin{lemma}[Inspect when the hypothesis holds]
In a problem where one is given a condition and asked to prove a result from it, first examine when the condition is satisfied, rather than trying to prove the result in a general case.
\end{lemma}
\paragraph{Application.}
If $p$ and $p^2+2$ are primes, then $p^3+2$ is also prime.
The prime $p$ must be odd. If $p>3$, then $p^2\equiv1\pmod3$, so $p^2+2$ is not prime.
Thus $p=3$.

Prove that the only integer solution of $x^2+y^2=x^2y^2$ is $x=y=0$.
\paragraph{Proof 1.}
$1=x^2y^2-x^2-y^2+1=(x^2-1)(y^2-1)$.
Thus $x^2-1=\pm1$ and $y^2-1=\pm1$, implying $x=y=0$.
\paragraph{Proof 3.}
From $x^2+y^2=x^2y^2$ the notes obtain $x=2x_1$ and $y=2y_1$, and then
\[
4x_1^2+4y_1^2=16x_1^2y_1^2
\Longrightarrow x_1^2+y_1^2=(2x_1y_1)^2
\Longrightarrow x_1=2x_2,\quad y_1=2y_2.
\]
% Source page 5 - right
Thus $x=2x_1=4x_2=\cdots=2^nx_n=\cdots$, so $(x_n)$ is strictly decreasing in magnitude.
Fermat's infinite descent principle gives the conclusion.

\setcounter{theorem}{17}\begin{lemma}[Pell's equation]
For $x^2-Dy^2=1$, if $(x_1,y_1)$ is a solution, then
\[
x_{n+1}=x_1x_n+Dy_1y_n,\qquad y_{n+1}=y_1x_n+x_1y_n.
\]
Equivalently,
\[
x_n=\frac12\bigl[(x_1+y_1\sqrt D)^n+(x_1-y_1\sqrt D)^n\bigr],\quad
y_n=\frac1{2\sqrt D}\bigl[(x_1+y_1\sqrt D)^n-(x_1-y_1\sqrt D)^n\bigr].
\]
\end{lemma}
\paragraph{Application: Romanian International Mathematical Olympiad (2000).}
Show that $x^3+y^3+z^2=t^4$ has infinitely many positive integer solutions.
The notes use
$(1+2^3+3^3+\cdots+n^3)=\bigl(n(n+1)/2\bigr)^2$, and hence
\[
\left(\frac{(n-2)(n-1)}2\right)^2+(n-1)^3+n^3
=\left(\frac{n(n+1)}2\right)^2.
\]
Requiring $n(n+1)/2=x^2$ gives
$n^2+n=2x^2$, or $(2n+1)^2-2(2x)^2=1$.
This is Pell's equation, with the recorded fundamental solution $(n,x)=(1,1)$.

\setcounter{theorem}{18}\begin{lemma}[Fermat's infinite descent principle]
There is no strictly decreasing sequence of positive integers.
This idea seems elementary but has many beautiful consequences.
\end{lemma}

% Source page 6 - left
\setcounter{theorem}{19}\begin{lemma}[Choose a modulus suited to the power]
Using an appropriate modulus for perfect powers often yields elegant solutions:
for cubes, modulo $7$ or $9$; for fourth powers, modulo $5$ or $16$;
for fifth powers, modulo $5$ or $11$. For instance, $x^5\equiv0,\pm1\pmod{11}$, and so on.
\end{lemma}

% Source page 6 - right
\setcounter{theorem}{20}\begin{lemma}[Factorization]
Some problems are resolved using simple factorization techniques; it is important to acquire skill in factorization.
\end{lemma}
% BLOG-CONTENT-END
\end{document}
